\section{Discussion}\label{discussion}

\subsection{When global queries resolve additional outcomes}\label{when-global-queries-resolve-additional-outcomes}

Ambiguity about which request produced an action matters to an audit if it can change the verdict on the target. Both local propagation and our method first check whether the full graph admits a feasible assignment. In the synchronized windows of the formal native cohort, local propagation assigned a bundle with a single candidate source, consumed that request's capacity, and determined the remaining assignments. The two methods therefore agreed. In the separate synchronized pilot window, two successful bundles had interchangeable sources. Local propagation could pin neither bundle to a target. The global query asks whether all bundles can still be assigned if the target receives no successful action. When the remaining eligible requests lack capacity, each feasible assignment gives the target a successful action even though the identity of its bundle remains ambiguous. Figure~\ref{fig:3} contrasts the two native structures. Figure~\ref{fig:9} and Appendix C show the corresponding capacity deficits in the sealed instances.

The ambiguity that remains after local propagation indicates where global queries can add a verdict. If the observations permit explanations both with and without a successful action for the target, the global verdict remains unresolved. If multiple assignments remain feasible but capacity already fixes the target outcome, a global query can return a determinate verdict. We observed this structure in the separate pilot window and included it in the designed synthetic instances. Its frequency in operational deployments remains to be measured.

\subsection{Evidence requirements and scope}\label{evidence-requirements-and-scope}

If the request records are incomplete and no verified link identifies a unique source, an unobserved request can still account for an action. Without a verified empty rule set before the write, a pre-existing rule can still explain reachability. If another request could have successfully written to the same resource during the window, the method leaves attribution of new admission unresolved. Paired evidence removal shows what additional evidence to collect. The competing-write check shows which other successful writes remain possible. These checks support separate verdicts for operation success and new admission.

The native experiments used a fixed build, a single resource, and one AC configuration. Six of 12 windows qualified in the initial cohort, whose inclusion criteria specified the expected post-state. In the formal cohort, 12 of 15 windows qualified, and collection eligibility and observed effects were recorded separately. Three formal C4 windows lacked end records under the six-second collection limit. We report the rechecks with a ten-second wait separately. We observed exchangeable source assignments in one separate pilot window. The same team generated and sealed the later synthetic instances using prespecified graph structures to test differences between methods under those structures.

\textbf{Capacity and multiple resources.} Suppose each bundle has a single initiating request and source code and configuration establish a trusted integer capacity \(k_v\) for each request \(v\). For each \(v\in Q\), the model can impose \(\sum_{a\in A}\mathbf{1}[f(a)=v]+h\mathbf{1}[v=q]\leq k_v\). Each dedicated external source retains unit capacity. Feasibility can be tested by capacitated matching or by expanding request \(v\) into \(k_v\) unit-capacity slots. A target query must consider all slots for the request, and forcing or allowing an unrecorded action must consume a slot. When request capacity exceeds one, reserving one slot requires changing line 10 of Algorithm 1, which removes all edges to the target in the unit-capacity case. Multiple resources also require checking the lifecycle and competing writes separately for each resource. These extensions require changes to the Boolean capacity field, the input contract, and the implementation, followed by validation. For asynchronous retries, the actual writer determines how to define the window and bundle. The applicable protocol or deployment policy determines when retransmission or renewal is allowed and which roles may initiate these operations.

\textbf{Temporal association.} The current model compares each request's processing interval with the operation's creation time within a single Server process instance. To associate evidence across nodes, first bound clock offset and drift over the observation window. The candidate intervals can then be widened to retain the actual source. Without these bounds, the temporal filter must be omitted and the corresponding explanations retained. Larger candidate sets can reduce the number of determinate verdicts.

\subsection{Interpreting collection and computation costs}\label{interpreting-collection-and-computation-costs}

Evidence collection must satisfy the verdict criteria within the service's cost constraints. Packet and event records support checks on the source of actions and the completeness of observations. Rule and fresh-connection records support verification of new admission. Each median paired difference describes a particular change in collection settings. The medians cannot be added across contrasts. Both paired knock-latency intervals for the addition of event logging include zero, leaving its separate contribution uncertain. The event writer encodes and writes its records as JSON under a mutex. Knock timing stops before synchronization at the end of collection. The individual contributions of scheduling, event writing, and packet capture still need to be measured.

Offline cost depends on the solver and how it is invoked. On the small graphs in the prospective sealed suite, median solver time was below one millisecond. Most fresh-process wall time was spent outside the solver. Preparation in Figure~\ref{fig:13} includes loading the method and processing its input and should be analyzed separately from solving. The benefit of reusing loaded modules and a candidate graph in a persistent process remains to be measured.

\section{Conclusion}\label{conclusion}

Temporary admission limits resource exposure to tens of seconds. Once access is observed, identifying the knock that opened it requires relating records from different parts of the path. The relevant records describe requests, rule state, and connection establishment. We formulate attribution over all explanations consistent with these observations. The evidence can establish that the target has a successful action even when several bundles could account for it. The method represents possible associations in a candidate graph. At the action level, queries that forbid or force edges distinguish necessary from possible success, as established by Propositions 1 and 2. With unit request capacity, feasible assignment reduces to bipartite matching. This lets global constraints resolve outcomes left open by local propagation. At the admission level, the method checks resource state before and after a write, verifies access over a fresh independent connection, and accounts for competing writes. It thereby reports operation success separately from the establishment of new admission.

On the native OpenNHP path, the evaluation distinguished entry into application processing, a successful rule write, and newly established admission. In the sealed synthetic comparisons, global verdicts agreed with exact enumeration for every instance. Global queries added determinate verdicts when capacity ruled out a feasible assignment without target success. Withdrawing request coverage made previously determinate verdicts unresolved. Full collection added a median of 0.620 ms to knock latency across paired blocks. Solving the small prospective sealed graphs took less than one millisecond at the median. The evaluation used one resource, at most one bundle per request, and one AC. The network cost measurements were taken under light load in a single virtual machine. Future work includes coupling capacity across resources, constructing candidates under bounded cross-node clock error, and determining when verdicts stabilize as uncertain observations are resolved.

\section*{Data availability}\label{data-availability}

The authors have archived the inputs, raw records, individual verdicts, and analysis scripts supporting these results. The common entry point, \path{artifacts/manifest.json}, lists each artifact's original path, purpose, and SHA-256 digest under its identifier. Appendices B--E identify the records needed for individual checks. The index is stored separately from the frozen originals, preserving cohort membership and the original digests.
